\documentclass[10pt,a4paper]{amsart}
\usepackage[margin=19mm]{geometry}
\usepackage{amsmath,amssymb,mathtools}
\usepackage[scr=rsfs,cal=euler]{mathalfa}
\usepackage{xfrac}
\usepackage[unicode,hidelinks,bookmarks=false]{hyperref}
\newcommand{\ie}{i.e.\ }
\newcommand{\cf}{cf.\ }
\newcommand{\resp}{respectively\ }
\newcommand{\Subsection}{\subsection}
\providecommand{\nobreakdash}{\nobreak\hbox{-}}
\setlength{\emergencystretch}{2em}
\multlinegap0pt
\usepackage{amssymb}
\usepackage{tikz}
\newtheorem{lem}{Lemma}
\def\bb{\mathbb}
\title{Ultrapowers of spectral subspaces}
\author{Hiroshi Ando, Isaac Goldbring}
\date{Source: arXiv:2605.21873v1 (CC BY 4.0)}
\begin{document}
\maketitle

\noindent\emph{Source and licence.} Ultrapowers of spectral subspaces. Version arXiv:2605.21873v1 (CC BY 4.0), distributed by arXiv under the Creative Commons Attribution 4.0 International licence, http://creativecommons.org/licenses/by/4.0/, as recorded in the arXiv licence metadata for this exact version and shown on the abstract page https://arxiv.org/abs/2605.21873v1. Full licence notice: \texttt{licenses/arxiv-2605.21873-CC-BY-4.0.txt}.\par\smallskip

\noindent\textit{Editorial scope.} Counted unit: the article's lem seq (source0.tex lines 345--361). The article's complete original proof of it is delivered with it (source0.tex lines 363--404, one \texttt{proof} environment of 42 lines). No mathematics is written, repaired or re-derived: the statement and the proof are the article's own lines, in the article's own notation, and no retained line is substituted or re-flowed.  No further source-local context is needed: every cross-reference of every retained line resolves inside the two delivered spans.  Nothing was omitted from any retained span, so the package carries no omission record.  No retained line contains a citation, so the wrapper carries no bibliography and no bibliography entry.  No retained line contains a figure environment, an image-inclusion command or drawing code, so no image is delivered and the package declares no asset dependency.  Editorial formatting only, in addition: an AMS \texttt{amsart} preamble that restates the article's own declarations and macros used by the retained lines, namely the environments \texttt{lem} on separate counters, together with the standard packages those lines need.  The retained environments print in this document as Lemma 1 in order of appearance, whereas the article prints the same items with the numbering of its own full document.  The complete native source of the article as distributed in the arXiv e-print for this exact version is bundled unchanged as \texttt{sources/201095/source0.tex}.
\begin{lem}\label{lem boundary seq}
Suppose \(F\subseteq \mathbb R\) is closed and \(a\in \partial F\). Then at least one of the following holds:
\begin{enumerate}
\item[{\rm (1)}] there exists a strictly increasing sequence \((a_n)_{n\in \bb N}\) such that
\begin{itemize}
    \item $a_n<a$ for all $n\in \bb N$,
    \item $a_n\notin F$ for all $n\in \bb N$, and 
    \item $\lim a_n=a$
\end{itemize}
\item[{\rm (2)}] there exists a strictly decreasing sequence \((a_n)_{n\in \bb N}\) such that
\begin{itemize}
    \item $a_n>a$ for all $n\in \bb N$,
    \item $a_n\notin F$ for all $n\in \bb N$, and 
    \item $\lim a_n=a$
\end{itemize}
\end{enumerate}
\end{lem}

\begin{proof}
Assume, towards a contradiction, that neither \((1)\) nor \((2)\) holds.  Since (1) fails, we claim that then there exists \(\varepsilon>0\) such that
\[
(a-\varepsilon,a)\subseteq F.
\]
Indeed, if no such \(\varepsilon\) existed, then for every \(n\in \mathbb N\) we could choose a point
\[
x_n\in (a-1/n,a)\setminus F.
\]
From these points we construct a strictly increasing subsequence converging to \(a\). Choose \(n_1\in \mathbb N\) arbitrarily. Having chosen \(n_k\), choose \(n_{k+1}>n_k\) large enough so that
\[
a-\tfrac1{n_{k+1}} > x_{n_k}.
\]
Then pick
\[
x_{n_{k+1}}\in (a-1/n_{k+1},a)\setminus F.
\]
Because \(x_{n_{k+1}}>a-\tfrac1{n_{k+1}}>x_{n_k}\), the sequence \((x_{n_k})\) is strictly increasing, each term lies outside \(F\), and \(x_{n_k}\to a\) since
\[
a-\tfrac1{n_k}<x_{n_k}<a.
\]
This contradicts the assumption that \((1)\) fails. Hence the claim is proved.

By the same argument, if \((2)\) fails, then there exists \(\delta>0\) such that
\[
(a,a+\delta)\subseteq F.
\]

Therefore, if both \((1)\) and \((2)\) fail, there exist \(\varepsilon,\delta>0\) such that
\[
(a-\varepsilon,a)\subseteq F
\qquad\text{and}\qquad
(a,a+\delta)\subseteq F.
\]
Since also \(a\in \partial F\subseteq \overline F=F\), this implies
\[
(a-\varepsilon,a+\delta)\subseteq F,
\]
so \(a\) is an interior point of \(F\). However, this contradicts \(a\in \partial F\).

Thus at least one of \((1)\) or \((2)\) must hold.
\end{proof}

\end{document}
